\documentclass[aps,prd,onecolumn,nofootinbib,superscriptaddress]{revtex4-2}

\usepackage[T1]{fontenc}
\usepackage[utf8]{inputenc}
\usepackage{amsmath,amssymb,amsthm,mathtools}
\usepackage{hyperref}
\usepackage{bm}
\usepackage{graphicx}
\usepackage{booktabs}

\newtheorem{theorem}{Theorem}
\newtheorem{proposition}{Proposition}
\newtheorem{lemma}{Lemma}
\newtheorem{corollary}{Corollary}
\newtheorem{remark}{Remark}

\begin{document}

\title{Wormhole Monopole Suppression: Per-End Conservation Laws,\\
the Harmonic Flux Sector, and Radiation Memory}

\author{Sergei Slobodov}
\email{sergei@slobodov.com}

\date{\today}

\begin{abstract}
On any asymptotically flat spacetime with multiple ends, the asymptotic monopole
data at each end---total electric charge $Q_a$ and ADM mass $M_a^{\mathrm{ADM}}$---are
separately conserved under all evolutions in which no charge or matter crosses
any asymptotic boundary.  For Maxwell, the proof is an immediate consequence of
Amp\`ere--Maxwell integrated over an asymptotic two-sphere and requires no
assumption about interior topology, throat geometry, traversability, or matter
content.  For gravity the corresponding statement follows from the multi-end
Regge--Teitelboim boundary-term analysis combined with the hypersurface
deformation algebra; it requires the standard admissibility of end-selecting
lapse functions with distinct asymptotic constants at each end, a condition
stated explicitly in Theorem~2.  These conservation laws apply to both
two-ended spacetimes (two distinct asymptotic regions, topology $\mathbb{R}\times
S^2$) and to single-ended spacetimes with a wormhole handle (one asymptotic
region, topology $\mathbb{R}^3$ with a handle); the harmonic-flux sector that
controls per-end bookkeeping lives in $H^2(\Sigma;\mathbb{R})$ for the
two-ended case and in $H^1(\Sigma;\mathbb{R})$ for the single-ended case, but
in both cases it is frozen by initial conditions and cannot be driven by local
source motion.

As an immediate corollary, the time-varying opposite-end monopole signal on
which the Dai--Stojkovic observational programme rests is forbidden at the
level of principle for both topologies: no interior source motion, throat
transit, or field configuration consistent with the hypotheses can produce a
time-varying $1/r$ signal at a distinct asymptotic end.  We diagnose precisely
where the Dai--Stojkovic inference breaks down: their static one-parameter
family is correctly computed, but it parametrises a slice through
(source-position, harmonic-sector) space that no dynamical trajectory traces.

We then identify and quantify the one effect that genuinely does exist and
partially validates the underlying intuition: radiation transiting the wormhole
throat---electromagnetic radiation from a passing star, or gravitational waves
from a compact-object merger in the vicinity of one mouth---backreacts at
second order on the throat geometry through the Isaacson (for gravitational
waves) or Poynting (for electromagnetic radiation) stress-energy.  By
per-end ADM conservation, any energy $E$ that crosses the throat from end $B$
to end $A$ permanently reduces the throat's contribution to the gravitational
flux by $E$; this is a monopole-type imprint of definite sign and can be
regarded as a throat memory effect.  Crucially, the sign is opposite to the
Dai--Stojkovic prediction (the throat becomes lighter, not heavier), the
magnitude is suppressed by two powers of the field amplitude relative to the
proposed signal, and the effect is undetectable with any foreseeable
instrument.  A model-dependent evanescent-bound argument additionally
suppresses all higher-multipole channels through long thin throats.  We conclude
that the specific monopole-channel signal on which several recent observational
proposals rest is incompatible with per-end conservation at the level of
principle, and that the surviving effects are both sign-reversed and far weaker
than claimed.
\end{abstract}

\maketitle

%% ============================================================
\section{Introduction}
\label{sec:intro}
%% ============================================================

The picture underlying several recent proposals for wormhole
detection~\cite{DaiStojkovic2019Observing,DaiStojkovic2020ReplyObserving,%
SimonettiEtAl2021SensitiveSearches,BaoHouZhang2023GWScattering,%
BambiStojkovic2021AstrophysicalWormholes} is intuitively appealing.
Consider a traversable wormhole whose two mouths both lie in the observable
universe.  A star or compact object on one side of the throat exerts a
gravitational (or electromagnetic) pull on the opposite side through the
wormhole connection.  As the source moves, the static exterior solution on the
opposite side changes: the coefficient of the $1/r$ term in the potential
varies with source position.  If this variation is large enough, a star
orbiting one mouth---say S2 orbiting Sgr~A*---could be seen to deviate from a
simple Keplerian orbit in a way attributable to an unseen mass on the other
side of the wormhole.  The same picture motivates proposals involving
gravitational wave scattering~\cite{BaoHouZhang2023GWScattering} and pulsar
timing~\cite{SimonettiEtAl2021SensitiveSearches}.

The main point of this paper is that the inference fails at the level of a
simple conservation law.  The coefficient of the $1/r$ term in the potential
at each asymptotic end---the ADM mass $M_a^{\mathrm{ADM}}$ or the total
electric charge $Q_a$---is a Hamiltonian boundary charge that is exactly
conserved under all evolutions preserving the asymptotic structure.  Source
motion in the interior, however complex, cannot modulate this boundary charge.
The time-varying far-zone monopole on which the cited proposals depend is
therefore absent.

This statement is not new in spirit.  Krasnikov's electrostatic
analysis~\cite{Krasnikov2008ElectrostaticInteraction,Krasnikov2020CommentObserving}
already isolates a fixed source-free wormhole flux parameter and notes that it
does not change under quasistatic source motion.  What we add is: (i) a clean
reformulation of the Maxwell result directly in terms of the per-end asymptotic
charge, proved using only Amp\`ere--Maxwell with no assumption about interior
structure; (ii) the corresponding gravitational statement, proved within the
multi-end Regge--Teitelboim framework and stated with its one remaining
technical condition made explicit; (iii) a treatment of both the two-ended
(two distinct asymptotic regions) and single-ended (one asymptotic region with
a wormhole handle) cases, clarifying which conservation law applies to which
topology and showing both rule out the claimed signal; (iv) an explicit
diagnosis of the Dai--Stojkovic family as a slice through (source position,
harmonic sector) parameter space that no evolution traces; (v) the
identification and quantification of the \emph{correct} subleading effect---a
radiation memory imprint on the throat geometry of definite sign and second
order in field amplitude.

\subsection*{Two topologies, one observational proposal}

A point requiring care from the outset is the relationship between the
mathematical setup and the observational framing.  The Dai--Stojkovic
construction uses two copies of $\mathbb{R}^3$ with balls removed and
boundaries identified---the mathematical structure of a two-ended wormhole with
a $\mathbb{R}\times S^2$ spatial slice and two distinct asymptotic
ends~\cite{DaiStojkovic2019Observing}.  Their observational proposal, however,
is phrased for ``our universe'': both wormhole mouths sit in a single connected
space, with the S2 star orbiting the Sgr~A* mouth~\cite{DaiStojkovic2019Observing,SimonettiEtAl2021SensitiveSearches}.
This is the \emph{single-ended} topology: one asymptotically flat region with a
handle connecting two interior 2-spheres.

The distinction matters because the conservation laws work differently in the
two cases.  For the two-ended case, each end has its own conserved ADM mass
$M_\pm^{\mathrm{ADM}}$, and the theorem directly forbids any change in either.
For the single-ended case, there is only one total $M^{\mathrm{ADM}}$, which is
trivially conserved, but the relevant observable---the apparent gravitational
mass near one mouth---is controlled instead by the harmonic-flux invariant
$Q_{\mathrm{wh}}$ living in $H^1(\Sigma;\mathbb{R})$, which is separately
conserved.  Both cases are closed by our analysis; the mechanisms differ
slightly, and both are described explicitly below.

\subsection*{The radiation memory effect}

Beyond the no-go result, we identify a genuine second-order effect that
partially vindicates the intuition behind the Dai--Stojkovic picture.
Radiation traversing the wormhole throat---light from a star, gravitational
waves from a black hole merger near one mouth---carries energy $E$ from the $B$
side to the $A$ side.  By the conservation of $M^{\mathrm{ADM}}$ at the $A$
end, the throat's gravitational flux contribution must decrease by exactly $E$
to compensate.  This is a permanent change: it persists after the radiation has
dispersed to infinity.  It constitutes a true monopole-type imprint on the
effective gravitational mass near the $A$ mouth.

This effect is real.  It has, however, three properties that distinguish it
sharply from the Dai--Stojkovic signal: (i) the sign is \emph{opposite}---the
throat becomes lighter, not heavier; (ii) the magnitude is second order in the
field amplitude ($O(h^2)$ for gravitational waves, $O(G S/c^4)$ for
electromagnetic radiation), compared to the first-order Dai--Stojkovic claim;
(iii) for any realistic source configuration the effect is many orders of
magnitude below current observational thresholds.  We estimate it for a
dark compact-object merger in the vicinity of the wormhole, which represents
the most favourable case.

\subsection*{Structure of the paper}

Section~\ref{sec:setup} defines the two relevant spacetime topologies and the
per-end observables.  Section~\ref{sec:topology} characterises the harmonic
flux sector for both topologies: $H^2$ for the two-ended case and $H^1$ for
the single-ended case.  Sections~\ref{sec:maxwell} and \ref{sec:adm} prove the
per-end conservation theorems for Maxwell and ADM mass.
Section~\ref{sec:synthesis} combines the topology and conservation results into
the dynamical reachability statement.  Section~\ref{sec:ds-diagnosis}
diagnoses the Dai--Stojkovic construction explicitly.
Section~\ref{sec:memory} develops the radiation memory effect.
Section~\ref{sec:topology-comparison} compares the one-ended and two-ended
cases systematically.  Section~\ref{sec:multiend} extends the framework to
$N$-ended spacetimes.  Section~\ref{sec:attenuation} establishes the
evanescent multipole attenuation bound.  Section~\ref{sec:phenomenology}
estimates all effects for the Sgr~A* scenario.
Section~\ref{sec:quasilocal} addresses the quasi-local mass objection.
Section~\ref{sec:discussion} concludes.  Three appendices supply the
cohomological details, the explicit transit accounting, and the Israel
junction-condition calculation for the throat memory.

%% ============================================================
\section{Setup and Definitions}
\label{sec:setup}
%% ============================================================

\subsection{Two-ended spacetimes}

Let $\mathcal{M}$ be a $(3+1)$-dimensional spacetime foliated by Cauchy slices
$\Sigma_t$, each a smooth oriented asymptotically flat 3-manifold with
\emph{two} ends $E_+$ and $E_-$.  The standard model is
\begin{equation}
  \Sigma \;\cong\; \mathbb{R}\times S^2,
\end{equation}
with each end corresponding to one of the two $\mathbb{R}$ semi-infinite rays.
The Dai--Stojkovic cut-and-paste construction---two copies of $\mathbb{R}^3$
with a ball removed from each, boundaries identified---is diffeomorphic to this
standard model.  We impose standard asymptotically Cartesian fall-off at each
end:
\begin{enumerate}
  \item $h_{ij} \equiv g_{ij}-\delta_{ij} = O(r_\pm^{-1})$ with appropriate
    fall-off of derivatives, at each end $E_\pm$.
  \item All charge and mass-energy sources are confined to a fixed compact
    interior region $K\subset\Sigma_t$, uniformly in $t$.
  \item No charge or matter crosses either asymptotic boundary
    $S_\pm(\infty)$ during the evolution.
  \item Electromagnetic and gravitational fields fall off at the rates required
    for finite asymptotic charges.\cite{ArnowittDeserMisner1962Dynamics,ReggeTeitelboim1974SurfaceIntegrals}
\end{enumerate}

For each end, define the asymptotic electric charge and ADM mass:
\begin{align}
  Q_\pm(t) &\;\equiv\;
    \frac{1}{4\pi}\lim_{R\to\infty}\int_{S_\pm(R)}\star F(t),
    \label{eq:Qpm-def}\\
  M_\pm^{\mathrm{ADM}}(t) &\;\equiv\;
    \frac{1}{16\pi G}\lim_{R\to\infty}
    \oint_{S_\pm(R)}\!\big(\partial_j h_{ij}-\partial_i h_{jj}\big)\,dS_i,
    \label{eq:Mpm-def}
\end{align}
with orientations chosen so that each counts the outward flux from end $E_\pm$.

\subsection{Single-ended spacetimes}

The observationally relevant case---both wormhole mouths in the same universe---is
the \emph{single-ended} topology.  The spatial slice is
\begin{equation}
  \Sigma_{\mathrm{se}} \;\cong\; \mathbb{R}^3\;\#\;(S^2\times S^1)
  \;\simeq\; \mathbb{R}^3\;\text{with a handle},
\end{equation}
constructed explicitly by removing two disjoint open balls $B_A,B_B\subset\mathbb{R}^3$
and identifying their boundaries $\partial B_A\cong\partial B_B\cong S^2$ via
a diffeomorphism.  This manifold has:
\begin{itemize}
  \item One asymptotically flat end $E_\infty$, with a single total charge
    $Q^{\mathrm{ADM}}$ and mass $M^{\mathrm{ADM}}$.
  \item Fundamental group $\pi_1(\Sigma_{\mathrm{se}})\cong\mathbb{Z}$
    (a loop threading the handle is non-contractible).
  \item First cohomology $H^1(\Sigma_{\mathrm{se}};\mathbb{R})\cong\mathbb{R}$,
    generated by a harmonic 1-form whose integral around the wormhole loop
    equals unity.
  \item Second cohomology $H^2(\Sigma_{\mathrm{se}};\mathbb{R})\cong\mathbb{R}$,
    generated by the area 2-form of the throat $S^2$.
\end{itemize}

The single global $M^{\mathrm{ADM}}$ is trivially conserved (it equals the
total energy of the universe).  The non-trivial conserved per-mouth quantity is
the harmonic-flux invariant $Q_{\mathrm{wh}}$ residing in $H^1$, described in
Sec.~\ref{sec:topology}.

\subsection{Dai--Stojkovic topology in context}

The Dai--Stojkovic paper~\cite{DaiStojkovic2019Observing} constructs its
wormhole as two copies of $\mathbb{R}^3$ with balls removed and boundaries
identified.  This is the \emph{two-ended} topology ($\Sigma\cong\mathbb{R}\times
S^2$, two AF ends), not the single-ended handle topology.  Their observational
proposal, however, places both mouths in our universe (the S2 star orbits the
Sgr~A* mouth), which corresponds to the single-ended framing.  We treat both
cases in the analysis below; the conservation laws that rule out the claimed
signal are present in both, and we note explicitly which applies to which
framing at each step.

%% ============================================================
\section{Wormhole Topology and the Harmonic Flux Sector}
\label{sec:topology}
%% ============================================================

The wormhole-specific physics enters at exactly one place: the cohomology of
the spatial slice determines a one-dimensional space of source-free field
configurations, labelled by a single real parameter $Q_{\mathrm{wh}}$.  The
topology of this parameter space differs between the one-ended and two-ended
cases, but in both cases $Q_{\mathrm{wh}}$ is frozen by initial conditions.

\subsection{Two-ended case: \texorpdfstring{$H^2$}{H2} and the throat flux}

\begin{lemma}[Standard two-ended slice]
\label{lem:standard-topology}
Let $\Sigma\cong\mathbb{R}\times S^2$ with two asymptotically flat ends.  Then
$H^2(\Sigma;\mathbb{R})\cong\mathbb{R}$, and the two end-spheres represent
opposite generators:
\begin{equation}
  [S_+] = -[S_-]\in H_2(\Sigma;\mathbb{R}).
\end{equation}
\end{lemma}

\begin{proof}
$\mathbb{R}$ is contractible, so $\Sigma\simeq S^2$, giving
$H^2(\Sigma;\mathbb{R})\cong H^2(S^2;\mathbb{R})\cong\mathbb{R}$.  A slab
$K=[-L,L]\times S^2$ has boundary $\{L\}\times S^2 - \{-L\}\times S^2$, so
the two caps are homologous.  The end spheres $S_\pm$ are large-$L$ deformations
with reversed relative orientation.
\end{proof}

\begin{corollary}[Single harmonic flux sector, two-ended]
\label{cor:single-sector-2e}
On a standard two-ended wormhole slice, the source-free part of any closed
Maxwell 2-form $\star F$ is characterised by a single real number
\begin{equation}
  Q_{\mathrm{wh}} \;\equiv\; \frac{1}{4\pi}\int_{S_{\mathrm{throat}}}\omega,
\end{equation}
where $\omega$ is the source-free part of $\star F$ and $S_{\mathrm{throat}}$
is any representative of the $H_2$ generator.  Equivalently, fixing a normalised
harmonic representative $h$ with $\frac{1}{4\pi}\int_{S_+(\infty)}h=+1$ and
$\frac{1}{4\pi}\int_{S_-(\infty)}h=-1$, every closed 2-form decomposes as
\begin{equation}
  \star F = \star F_{\mathrm{src}} + 4\pi Q_{\mathrm{wh}}\,h,
\end{equation}
with $\star F_{\mathrm{src}}$ source-attached and $Q_{\mathrm{wh}}\in\mathbb{R}$
labelling the topological sector.
\end{corollary}

The decisive feature of $Q_{\mathrm{wh}}$ is cohomological: it is the
coefficient of the unique closed non-exact 2-form on $\Sigma$.  This is the
original Wheelerian ``charge without charge.''~\cite{MisnerWheeler1957ClassicalPhysics,Wheeler1962Geometrodynamics}

\subsection{Single-ended case: \texorpdfstring{$H^1$}{H1} and the
circulation flux}

For the single-ended wormhole $\Sigma_{\mathrm{se}}$ with
$\pi_1\cong\mathbb{Z}$ and $H^1(\Sigma_{\mathrm{se}};\mathbb{R})\cong\mathbb{R}$,
the harmonic-flux sector is different in character.

\begin{lemma}[Harmonic flux sector, single-ended]
\label{lem:single-ended-sector}
On $\Sigma_{\mathrm{se}}$, the space of source-free static electric fields is
one-dimensional and generated by a harmonic 1-form $\alpha$ satisfying
$d\alpha=0$ and $d\star_3\alpha=0$, with
$\oint_{\gamma_{\mathrm{wh}}}\alpha = 1$ for the non-contractible loop
$\gamma_{\mathrm{wh}}$ threading the handle.  Any static electric field on
$\Sigma_{\mathrm{se}}$ decomposes as
\begin{equation}
  E^\flat = -d\phi + Q_{\mathrm{wh}}\,\alpha,
\end{equation}
where $\phi$ is a globally defined single-valued scalar potential (well-defined
because $H^1(\Sigma_{\mathrm{se}};\mathbb{R})\cong\mathbb{R}$ but the kernel of
the scalar Laplacian is trivial by the maximum principle) and $Q_{\mathrm{wh}}$
is the single free parameter labelling the harmonic-1-form sector.
\end{lemma}

The throat flux in this case is
\begin{equation}
  Q_{\mathrm{wh}} = \frac{1}{4\pi}\int_{S^2_{\mathrm{throat}}}\star_3 E^\flat
  = \frac{1}{4\pi}\int_{S^2_{\mathrm{throat}}}\star_3\alpha,
\end{equation}
which measures the net electric flux threading the wormhole handle.

\begin{remark}[Why scalar potential uniqueness does not imply $Q_{\mathrm{wh}}=0$]
The maximum principle gives uniqueness of $\phi$ (the scalar part), but the
harmonic 1-form $\alpha$ is not $d$ of any globally defined function: it
requires going around the non-contractible loop to contribute.  The free
parameter $Q_{\mathrm{wh}}$ lives in the kernel of the 1-form Laplacian, not
in the kernel of the scalar Laplacian.  The two kernels coincide only when
$H^1=0$, i.e., on simply connected spaces.
\end{remark}

\begin{remark}[Comparison of the two cases]
In the two-ended case, $Q_{\mathrm{wh}}$ lives in $H^2(\Sigma;\mathbb{R})$ and
is a property of the 2-form field strength $F$ on the spatial slice.  In the
single-ended case, it lives in $H^1(\Sigma_{\mathrm{se}};\mathbb{R})$ and is a
property of the 1-form electric field $E^\flat$.  Both are rigid under local
field-theoretic evolution, but the rigidity is enforced by a 4-dimensional
Stokes argument in both cases (Appendix~\ref{app:cohomology}).
\end{remark}

\subsection{Multi-end generalization}

Wormholes with $N>2$ asymptotically flat ends, or with additional independent
2-cycles, enlarge the harmonic sector: $H^2(\Sigma;\mathbb{R})$ becomes
$N-1$-dimensional for an $N$-ended spacetime with a single connecting handle
(one free flux parameter per pair of ends), and can be larger for more complex
topologies.  The conservation theorems of Sections~\ref{sec:maxwell} and
\ref{sec:adm} apply verbatim to each end separately.  We restrict to two-ended
and single-ended cases in the detailed treatment and address the $N$-ended
generalization in Sec.~\ref{sec:multiend}.

%% ============================================================
\section{Per-End Conservation: Maxwell Case}
\label{sec:maxwell}
%% ============================================================

\subsection{Two-ended case}

\begin{theorem}[Per-end charge conservation, two-ended]
\label{thm:Q-conservation-2e}
Let $(F(t),\rho(t),\bm{j}(t))$ be a smooth Maxwell evolution on the two-ended
slice $\Sigma\cong\mathbb{R}\times S^2$, with sources compactly supported in a
fixed interior region $K$ and $\bm{B}$ falling off fast enough that
$(\nabla\times\bm{B})$ has vanishing surface flux at $S_\pm(\infty)$.  Then
\begin{equation}
  \frac{dQ_\pm}{dt} = 0.
\end{equation}
\end{theorem}

\begin{proof}
Fix end $E_+$ and a sphere $S_+(R)$ exterior to $K$.  Integrate
Amp\`ere--Maxwell over $S_+(R)$:
\begin{equation}
  \frac{d}{dt}\oint_{S_+(R)}\bm{E}\cdot d\bm{A}
  = \oint_{S_+(R)}(\nabla\times\bm{B})\cdot d\bm{A}
  - 4\pi\oint_{S_+(R)}\bm{j}\cdot d\bm{A}.
\end{equation}
The current term vanishes because $\bm{j}$ is supported in $K$, disjoint from
$S_+(R)$.  The curl term vanishes by $\nabla\cdot(\nabla\times\bm{B})\equiv 0$
applied to the exterior region $\mathrm{ext}_+(R)$: the boundary of
$\mathrm{ext}_+(R)$ in $\Sigma$ consists of $S_+(R)$ and $S_+(\infty)$, and
the latter contributes zero by the asymptotic fall-off assumption on $\bm{B}$.
Hence $\frac{d}{dt}\oint_{S_+(R)}\bm{E}\cdot d\bm{A}=0$; taking $R\to\infty$
gives $dQ_+/dt=0$.
\end{proof}

\begin{remark}[What the proof does not use]
The proof uses only Amp\`ere--Maxwell, the divergence theorem in one
asymptotic end, and the fall-off of $\bm{B}$.  It does \emph{not} use the
topology of the interior, the existence of a throat, traversability, the
matter content, any junction conditions, or any symmetry.  Whatever happens in
the interior---including transit of charged bodies through any throat---is
invisible to a sphere lying outside $K$ at one end.
\end{remark}

\subsection{Single-ended case}

For the single-ended wormhole, there is one total charge $Q^{\mathrm{ADM}}$,
trivially conserved.  The per-mouth quantity is $Q_{\mathrm{wh}}$, the
harmonic-1-form sector coefficient.

\begin{theorem}[Per-mouth flux conservation, single-ended]
\label{thm:Q-conservation-1e}
Let $(F(t),\rho(t),\bm{j}(t))$ be a smooth Maxwell evolution on
$\Sigma_{\mathrm{se}}$, with sources compactly supported in a region $K$ that
misses some representative of the handle's non-contractible loop.  Then
\begin{equation}
  \frac{dQ_{\mathrm{wh}}}{dt} = 0.
\end{equation}
\end{theorem}

\begin{proof}
Choose a 2-cycle $S\subset\Sigma_{\mathrm{se}}$ representing the $H_2$
generator (a 2-sphere separating the throat from the rest of the manifold)
that is disjoint from $K$ on the time interval $[0,T]$.  The 3-chain
$S\times[0,T]$ in 4-dimensional spacetime has boundary
$S\times\{T\}-S\times\{0\}$.  Applying $d(\star_4 F)=4\pi\star_4 J$ and
integrating:
\begin{equation}
  \int_{S\times\{T\}}\star_4 F - \int_{S\times\{0\}}\star_4 F
  = 4\pi\int_{S\times[0,T]}\star_4 J = 0,
\end{equation}
since $S\times[0,T]$ is disjoint from $J$.  Hence $Q_{\mathrm{wh}}(T)=Q_{\mathrm{wh}}(0)$.
\end{proof}

\begin{corollary}[D-S monopole induction forbidden, single-ended Maxwell]
A source confined to the neighborhood of mouth $B$ cannot induce a
time-varying flux $Q_{\mathrm{wh}}$ through the throat, and hence cannot
produce a time-varying apparent monopole at mouth $A$.
\end{corollary}

\begin{remark}[Transit vs.\ in-end motion]
If the source \emph{traverses} the throat, every 2-sphere representative of
$H_2$ is intersected, the 3-chain $S\times[0,T]$ is no longer disjoint from
$J$, and $Q_{\mathrm{wh}}$ shifts by the transported charge.  The total
$Q^{\mathrm{ADM}}$ is still conserved, and the apparent charge near each mouth
adjusts: the mouth through which the charge exits gains apparent charge
$+q_{\mathrm{transit}}$, the other mouth loses it, and the combined bookkeeping
is consistent with total conservation.
\end{remark}

%% ============================================================
\section{Per-End Conservation: ADM Mass}
\label{sec:adm}
%% ============================================================

\subsection{Two-ended case}

The gravitational analogue holds under two explicit conditions.

\begin{theorem}[Per-end ADM conservation, two-ended]
\label{thm:M-conservation}
Under the assumptions of Sec.~\ref{sec:setup}, and assuming (i) that $\Sigma$
has no additional physical boundary components beyond the asymptotic ends, and
(ii) that for each end $E_a$ there exists a smooth lapse $N_a$ satisfying
$N_a\to 1$ at $E_a$, $N_a\to 0$ at every other end, and the Regge--Teitelboim
admissibility conditions at each end,\cite{ReggeTeitelboim1974SurfaceIntegrals,Chrusciel2020ConstraintEquations} the ADM masses are independent of $t$:
\begin{equation}
  \frac{dM_\pm^{\mathrm{ADM}}}{dt} = 0.
\end{equation}
\end{theorem}

\begin{proof}
The gravitational Hamiltonian is
\begin{equation}
  H[N,N^i] = \int_\Sigma(N\mathcal{H}+N^i\mathcal{H}_i)\,d^3x + B[N,N^i],
\end{equation}
where $B$ collects asymptotic boundary terms.  On the constraint surface the
bulk integral vanishes.

\textbf{Step 1.} For each end $E_a$, the end-selecting lapse $N_a$ (assumption
ii) gives $H_a:=H[N_a,0]$, which equals $M_a^{\mathrm{ADM}}$ on the constraint
surface.

\textbf{Step 2.} The hypersurface deformation algebra
(HDA)~\cite{HojmanKucharTeitelboim1976Geometrodynamics} gives an exact identity:
\begin{equation}
  \{H[N_a,0],H[N_b,0]\} = H[0,\beta_{ab}],\qquad
  \beta_{ab}^i = h^{ij}(N_a\partial_jN_b-N_b\partial_jN_a).
\end{equation}

\textbf{Step 3.} Since each $N_a\to\mathrm{const}$ at every end, $\partial_jN_a\to 0$
asymptotically, so $\beta_{ab}^i\to 0$ at every boundary.

\textbf{Step 4.} In the RT framework, $H[0,\beta_{ab}]$ vanishes weakly:
its boundary term is $\int_{\partial\Sigma}\beta_{ab}^i(\text{metric data})=0$
(since $\beta_{ab}\to 0$ at all ends), and the bulk integral vanishes on the
constraint surface.  Assumption (i) excludes additional internal boundary
terms.  Therefore $\{H_a,H_b\}\approx 0$ weakly.

\textbf{Step 5.} Any admissible evolution Hamiltonian takes the form
$H_{\mathrm{evol}}=\sum_b c_b H_b+H_{\mathrm{gauge}}$,
so
\begin{equation}
  \frac{dM_a^{\mathrm{ADM}}}{dt}
  = \{H_a,H_{\mathrm{evol}}\}
  = \sum_b c_b\{H_a,H_b\}+\{H_a,H_{\mathrm{gauge}}\}
  \approx 0.
\end{equation}
\end{proof}

\begin{remark}[ADM vs.\ Bondi mass]
\label{rem:adm-vs-bondi}
$M_\pm^{\mathrm{ADM}}$ at spatial infinity $i^0$ is exactly conserved.  The
Bondi mass at null infinity $\mathcal{I}^\pm$ decreases via outgoing
gravitational radiation:
\begin{equation}
  \frac{dM_\pm^{\mathrm{Bondi}}}{du}
  = -\frac{1}{4\pi G}\oint|\mathcal{N}|^2\,d\Omega,
\end{equation}
where $\mathcal{N}$ is the news function and $u$ retarded
time.\cite{BondiVanderBurgMetzner1962Gravitational,Sachs1962Asymptotic}
The ADM mass includes all energy that will ever radiate, which is why it is
conserved without a radiation correction.  Any argument that ``mass should
change because gravitational waves carried it away'' is an argument about
Bondi mass, not ADM, and even at the Bondi level it cannot \emph{add} mass to
an end where none was present.
\end{remark}

\subsection{Single-ended case}
\label{sec:adm-single-ended}

For $\Sigma_{\mathrm{se}}$ with one AF end, $M^{\mathrm{ADM}}$ is the total
energy of the spacetime and is trivially conserved.  The per-mouth gravitational
analogue of $Q_{\mathrm{wh}}$ is the gravitational throat flux
$Q_{\mathrm{wh}}^{\mathrm{grav}}$---roughly, the integral of the extrinsic
curvature contribution over the throat 2-sphere---which governs how much
gravitational flux from the $B$-side sources threads through to the $A$ side.

\textbf{Linearized gravity.}  In linearized GR around a fixed wormhole
background, the metric perturbation $h_{\mu\nu}$ obeys a linearized wave
equation.  The linearized Bianchi identity $\partial_\mu G^{\mu\nu}_{\text{lin}}=0$
gives a structure identical to Maxwell, and the same 4-dimensional Stokes
argument applies: for matter not transiting the throat,
$\partial_t Q_{\mathrm{wh}}^{\mathrm{grav}}=0$ to linear order.

\textbf{Nonlinear corrections.}  The full Einstein equations are nonlinear.
The gravitational field itself acts as an effective stress-energy source
(through the Landau--Lifshitz pseudo-tensor, or equivalently through the
quadratic terms in the Ricci tensor), and this effective source is nonzero even
in vacuum.  At second order in the metric perturbation, this can in principle
drive a change in $Q_{\mathrm{wh}}^{\mathrm{grav}}$ even without matter
transiting.  However, as we show in Sec.~\ref{sec:memory}, the only physical
process capable of doing so---gravitational radiation transiting the throat---
produces a throat mass \emph{deficit} of sign opposite to the Dai--Stojkovic
prediction, and at second order in field amplitude.  In no case does a
first-order monopole induction occur.

%% ============================================================
\section{Synthesis: Dynamical Reachability}
\label{sec:synthesis}
%% ============================================================

Combining the topological input of Sec.~\ref{sec:topology} with the
conservation laws of Secs.~\ref{sec:maxwell} and \ref{sec:adm} yields the
complete rigidity statement.

The space of static Maxwell solutions on a fixed two-ended wormhole is
parameterized by three inputs: the source distribution $\rho$, the per-end
charges $(Q_+,Q_-)$ with $Q_++Q_-=q_{\mathrm{total}}$, and asymptotic boundary
data.  For fixed $\rho$ and boundary data, the static family is
one-dimensional, equivalently parameterized by $Q_{\mathrm{wh}}\in H^2(\Sigma;\mathbb{R})\cong\mathbb{R}$.

Any continuous dynamical evolution starting from a given static configuration
keeps $Q_\pm$ fixed (Theorem~\ref{thm:Q-conservation-2e}).  The dynamical
end-state therefore has the same $(Q_+,Q_-)$ as the initial configuration,
with whatever source distribution the evolution has produced.  The static
end-state belongs to the same harmonic-flux sector as the initial state, not
to the sector selected by the Dai--Stojkovic matching conditions at the final
source position.

\begin{corollary}[Dynamical reachability]
\label{cor:reachability}
Two static Maxwell configurations on the same wormhole geometry with the same
total source charge but different per-end charges $(Q_+,Q_-)$ cannot be
connected by any dynamical evolution satisfying the hypotheses of
Theorem~\ref{thm:Q-conservation-2e}.  They occupy different harmonic-flux
sectors.
\end{corollary}

\begin{corollary}[Same for ADM, conditional]
\label{cor:reachability-adm}
Under the conditions of Theorem~\ref{thm:M-conservation}, two quasi-static
spacetime configurations with the same total mass-energy but different per-end
ADM masses are dynamically disconnected.
\end{corollary}

The asymmetry in these corollaries should be stressed.  They do not say that
two distinct static solutions cannot both exist---they manifestly do, parameterized
by the harmonic sector.  They say only that one cannot evolve from one to the
other under dynamics respecting the conservation laws.  The Dai--Stojkovic
inference error is precisely: existence of a static family $\Rightarrow$
existence of an evolution tracing that family.

Figure~\ref{fig:fieldlines} provides an illustrative example of the
bookkeeping in the Dai--Stojkovic short-throat model over six stages of
source transit.

\begin{figure*}[tb]
\centering
\includegraphics[width=0.95\textwidth]{fig1_fieldlines.pdf}
\caption{\textit{Illustrative example} of per-end conservation in the
Dai--Stojkovic short-throat model (full multipole sum, $\ell_{\max}=30$).
\textit{Geometry:} the gray stadium represents the wormhole throat; side $+$
(blue) occupies $u<-L/2$, side $-$ (red) occupies $u>+L/2$.  \textit{Colour:}
line colour encodes field magnitude on a shared logarithmic scale.  Panels
(a)--(c): source on side $+$ approaching the mouth.  Side $-$ carries only a
dipolar response; no $1/r$ monopole is present.  Panels (d)--(f): source has
emerged on side $-$.  Side $+$ now supports a monopole of strength $q$ centred
at the left mouth, as required by the harmonic-flux sector that conservation
pins at that end.  Per-end conservation (Theorem~\ref{thm:Q-conservation-2e})
is model-independent; the field pattern shown here is one realisation of what
it implies.}
\label{fig:fieldlines}
\end{figure*}

%% ============================================================
\section{Diagnosis of the Dai--Stojkovic Construction}
\label{sec:ds-diagnosis}
%% ============================================================

\subsection{What they compute}

The Dai--Stojkovic setup~\cite{DaiStojkovic2019Observing} is two copies of
flat $\mathbb{R}^3$, each with a ball of radius $R$ removed, boundaries
identified.  A point charge $q$ at distance $A$ on side 1 (``other space'') is
matched to a potential on side 2 (``our space'') via continuity conditions
$V_1(R)=V_2(R)$ and $\partial_{r_1}V_1|_R = -\partial_{r_2}V_2|_R$, giving
\begin{equation}
\label{eq:DS-Q}
Q_2 = \frac{qR}{2A}, \qquad Q_1 = q - \frac{qR}{2A}.
\end{equation}
This is a correct static computation.  For each $A$ it identifies a valid static
solution on the cut-and-paste manifold: a one-parameter family
$\{F_A\}_{A>R}$ of static solutions with $Q_2(A)=qR/(2A)$ varying with source
position.

\subsection{The D-S picture and why it is appealing}
\label{sec:ds-picture}

Figure~\ref{fig:ds-claimed} illustrates why the Dai--Stojkovic picture is
intuitive.  In six panels, a source moves from far on side 1 toward the throat
and through to side 2.  In each panel, the Dai--Stojkovic matching conditions
determine a monopole coefficient $Q_2(A)=qR/(2A)$ on side 2 that grows as the
source approaches the throat, peaks when the source is at the mouth, and then
decreases as the source recedes on side 2.  An observer orbiting side 2's mouth
would, in this picture, see a time-varying $1/r^2$ gravitational acceleration
whose amplitude follows $q R/(2A(t))$.

\begin{figure*}[tb]
\centering
\includegraphics[width=0.95\textwidth]{fig_ds_claimed.pdf}
\caption{The Dai--Stojkovic picture of wormhole monopole induction, as they
describe it.  Six panels show the source (solid dot) at successive positions on
side 1 (blue) approaching and traversing the throat (gray) to side 2 (red).
At each stage, the Dai--Stojkovic matching conditions produce a growing
monopole contribution $Q_2=qR/(2A)$ on side 2 (dashed field lines).  An
observer on side 2 would see a time-varying $1/r^2$ signal tracking the source
position.  This is the picture our conservation theorems rule out.}
\label{fig:ds-claimed}
\end{figure*}

\subsection{The dynamical inference that fails}

The breakdown is in the implicit step that treats the family $\{F_A\}$ as a
continuous dynamical trajectory.  Under this reading, $A$ plays the role of a
time parameter and an observer on side 2 sees $Q_2(t)=qR/(2A(t))$.

By Theorem~\ref{thm:Q-conservation-2e}, no such evolution exists.  As the
source moves from $A_0$ to $A_1$, $Q_2$ remains at its initial value
$Q_2(A_0)$.  The end-state is a static configuration with source at $A_1$ and
$Q_2$ still equal to $qR/(2A_0)$, not $qR/(2A_1)$.

\subsection{Where the harmonic-flux drift is hidden}

The two static solutions with source at $A_1$---the Dai--Stojkovic solution
$F_{A_1}^{\mathrm{DS}}$ (with $Q_2=qR/(2A_1)$) and the dynamical end-state
$F_{A_1}^{\mathrm{dyn}}$ (with $Q_2=qR/(2A_0)$)---differ by a source-free
closed 2-form:
\begin{equation}
\star F_{A_1}^{\mathrm{dyn}} - \star F_{A_1}^{\mathrm{DS}}
= 4\pi\,\Delta Q_{\mathrm{wh}}\,h, \qquad
\Delta Q_{\mathrm{wh}} = \frac{qR}{2}\!\left(\frac{1}{A_0}-\frac{1}{A_1}\right).
\end{equation}
Equivalently, tracing the Dai--Stojkovic curve requires the harmonic flux to
drift:
\begin{equation}
Q_{\mathrm{wh}}^{\mathrm{DS}}(A) = Q_{\mathrm{wh}}^{\mathrm{DS}}(A_0)
  - \frac{qR}{2}\!\left(\frac{1}{A_0}-\frac{1}{A}\right).
\end{equation}
The harmonic flux is a topological invariant; it cannot drift under any local
field-theoretic dynamics.  The Dai--Stojkovic family is therefore a family of
solutions to a sequence of differently-sectored static problems, not the
dynamical history of any physical wormhole.

Figure~\ref{fig:multipoles} makes this quantitative at the level of the
side-2 multipole coefficients.

\begin{figure}[t]
\centering
\includegraphics[width=0.88\textwidth]{fig2_multipoles.pdf}
\caption{Side-2 multipole coefficients $|B_\ell|$ versus source distance $A/R$
in the Dai--Stojkovic short-throat geometry.  Per-end conservation pins
$B_0\equiv 0$ (thick black); the Dai--Stojkovic prescription gives
$B_0(A)=qR/(2A)$ (dashed), which requires a forbidden harmonic-flux drift.
Higher multipoles $\ell=1,2,3,4$ (solid) are unconstrained and represent the
real cross-throat channels.}
\label{fig:multipoles}
\end{figure}

\subsection{Single-ended re-examination}

In the single-ended framing, there is one mouth $A$ (the observable end, near
Sgr~A*) and one mouth $B$ (the other mouth, somewhere in our universe).  The
total charge/mass is trivially conserved.  The D-S monopole signal would
require the \emph{apparent} monopole at mouth $A$---i.e., the gravitational
field at S2's orbital distance---to change as mass moves from the $B$-side
region to the $A$-side region.

This is controlled by $Q_{\mathrm{wh}}$ (for Maxwell, residing in $H^1$).
By Theorem~\ref{thm:Q-conservation-1e}, $Q_{\mathrm{wh}}$ is conserved as long
as the source does not transit the throat.  A source \emph{staying} on the $B$
side cannot change the flux threading the throat and hence cannot change the
apparent monopole at mouth $A$.  The claimed induction effect is absent.

For a source that \emph{does} transit from $B$ to $A$ (matter physically
arriving at Sgr~A* through the wormhole): the charge/mass does appear near
mouth $A$, and S2 does see a changed gravitational field.  But this is not
``induction''---it is the transited mass being physically present at mouth $A$.
For this to explain the Dai--Stojkovic signal, the matter would need to be
arriving continuously and undetected, which is a separate astrophysical
hypothesis and not a consequence of the wormhole topology.

%% ============================================================
\section{Radiation Memory on the Wormhole Throat}
\label{sec:memory}
%% ============================================================

\subsection{Physical mechanism}

The conservation laws established above apply to first-order (linear) processes:
source motion and transit.  At second order in the field amplitude, a genuine
physical effect exists that partially vindicates the Dai--Stojkovic intuition,
while differing from it in sign and magnitude.

When radiation---electromagnetic (light from a star) or gravitational (waves
from a compact-object merger)---transits the wormhole throat from mouth $B$ to
mouth $A$, it carries energy $E$ from the $B$ region to the $A$ region.
By the conservation of $M_+^{\mathrm{ADM}}$ (two-ended case, Theorem~\ref{thm:M-conservation}),
this energy redistribution must be compensated: the throat's gravitational flux
contribution $Q_{\mathrm{wh}}^{\mathrm{grav}}$ must decrease by $E$ to keep
$M_+^{\mathrm{ADM}}$ constant.  Since this change persists after the radiation
has dispersed to infinity, it is a \emph{permanent} change in the throat
geometry---a memory effect.

The compensation mechanism is the gravitational analogue of what happens in the
Maxwell case when a charge transits the throat (Section~\ref{sec:maxwell}):
whatever energy/charge exits through mouth $A$, the throat ``releases'' an
equal and opposite contribution so that the total at the $A$ end remains
constant.  The wormhole end always gains a contribution from the outgoing
radiation and loses the same amount from the throat flux, with zero net change
at $A$ infinity.

\subsection{Sign and order of the effect}

The key points are:

\begin{enumerate}
\item \textbf{Sign.} The throat mass contribution decreases: $\delta
M_{\mathrm{throat}} = -E < 0$.  An observer near mouth $A$ (but far from the
throat) sees a slightly \emph{lighter} throat after the radiation has passed.
This is opposite to the Dai--Stojkovic prediction of an added positive mass.

\item \textbf{Order.} The effect is second-order in the field amplitude.  For
gravitational waves, $E\sim h^2$ (Isaacson stress-energy); for electromagnetic
radiation, $E\sim S/c$ (Poynting flux) with the gravitational backreaction
proportional to $GE/c^4$.  In both cases the monopole imprint at mouth $A$
scales as $\delta M_{\mathrm{throat}}/M_{\mathrm{throat}} \sim E/M_{\mathrm{throat}} c^2$,
which for any realistic source is many orders of magnitude below the
Dai--Stojkovic first-order effect.

\item \textbf{Locality.}  The effect is localised near the throat.  For an
observer at distance $r$ from mouth $A$ with $r\gg R_{\mathrm{throat}}$, the
throat mass deficit appears as a monopole correction to the gravitational
potential:
\begin{equation}
  \delta\Phi(r) \approx -\frac{G\,\delta M_{\mathrm{throat}}}{r}
  = +\frac{G E}{r} > 0.
\end{equation}
(The potential becomes less negative as the throat loses mass; a distant observer
experiences a slightly weaker gravitational pull from mouth $A$.)
\end{enumerate}

\subsection{Memory interpretation}

The permanent character of the throat mass change distinguishes this from
transient oscillatory effects.  In the gravitational wave case, the analogous
phenomenon in asymptotically flat spacetime is the Christodoulou
memory~\cite{Christodoulou1991Memory}---a permanent DC shift in the metric induced by the nonlinear
stress-energy of the radiation.  The wormhole throat memory is a distinct effect:
it is a change in the $\ell=0$ (monopole) component of the throat's gravitational
mass, not a transverse-traceless (TT) displacement of test masses.  Both
effects are real and both are permanent; both are second-order in amplitude.

The wormhole throat memory is in this sense a new addition to the catalogue of
gravitational memory effects: a permanent monopole change in the wormhole
throat mass, sourced by radiation transit, of sign opposite to naive induction
intuition.

\begin{remark}[Relation to Christodoulou memory]
The Christodoulou memory~\cite{Christodoulou1991Memory} is a permanent
displacement of test masses due to the nonlinear stress-energy of outgoing
gravitational radiation; it is a transverse-traceless (TT) $\ell \ge 2$
effect measurable by interferometers.  The wormhole throat memory is
distinct: it is an $\ell = 0$ (monopole) shift in the throat's effective
gravitational mass, sourced by radiation transit through the handle.  Both
are real, permanent, and second-order in field amplitude.
\end{remark}

\subsection{Electromagnetic radiation: starlight}

Consider a star of luminosity $L_*$ at distance $d$ from wormhole mouth $B$,
whose light partially transits the throat to mouth $A$.  The solid angle
subtended by the throat (of radius $R_{\mathrm{throat}}$) from the star's
location is $\Omega = \pi R_{\mathrm{throat}}^2/d^2$.  The rate of energy
transit is
\begin{equation}
  \dot{E}_{\mathrm{EM}} = L_* \cdot \frac{\Omega}{4\pi}
  = \frac{L_* R_{\mathrm{throat}}^2}{4 d^2}.
\end{equation}
The resulting rate of throat mass change is
\begin{equation}
  \dot{M}_{\mathrm{throat}} = -\frac{\dot{E}_{\mathrm{EM}}}{c^2}
  = -\frac{L_* R_{\mathrm{throat}}^2}{4\, c^2\, d^2},
\end{equation}
which integrated over a time $\Delta t$ gives
\begin{equation}
  \delta M_{\mathrm{throat}}^{\mathrm{EM}}
  = -\frac{L_*\,R_{\mathrm{throat}}^2}{4\,c^2\,d^2}\,\Delta t.
\end{equation}
For a solar-type star ($L_*=L_\odot\approx 4\times 10^{26}$\,W) at
$d=1$\,AU from a macroscopic wormhole throat ($R_{\mathrm{throat}}=10^3$\,m,
a generous estimate for a stellar-mass traversable wormhole), accumulated over
one year ($\Delta t\approx 3\times 10^7$\,s):
\begin{equation}
  |\delta M_{\mathrm{throat}}^{\mathrm{EM}}|
  \approx
  \frac{4\times 10^{26}\times 10^6}{4\times (9\times 10^{16})\times (1.5\times 10^{11})^2}
  \times 3\times 10^7
  \;\mathrm{kg}
  \approx 10^{-11}\;\mathrm{kg}.
\end{equation}
This is sub-nanogram per year.  Compare with the wormhole's own
Schwarzschild mass $M_{\mathrm{throat}}\sim M_\odot\approx 2\times 10^{30}$\,kg:
\begin{equation}
  \frac{|\delta M_{\mathrm{throat}}^{\mathrm{EM}}|}{M_{\mathrm{throat}}}
  \;\lesssim\; 5\times 10^{-42}\quad\text{per year.}
\end{equation}

\subsection{Gravitational waves: dark compact-object merger}

The most favourable scenario for the gravitational wave memory effect is a
compact binary merger (black hole--black hole or neutron star--black hole) in
the vicinity of wormhole mouth $B$.  Such a merger radiates a fraction
$\epsilon\sim 0.05$ of its total mass as gravitational waves, so
$E_{\mathrm{GW}}\approx\epsilon M_{\mathrm{merger}}c^2$.  Of this, the
fraction threading the throat is the solid angle factor
$\Omega_{\mathrm{throat}}/(4\pi) = R_{\mathrm{throat}}^2/(4 D^2)$, where $D$
is the merger-to-throat distance.

\begin{equation}
  \delta M_{\mathrm{throat}}^{\mathrm{GW}}
  = -\frac{\epsilon M_{\mathrm{merger}} R_{\mathrm{throat}}^2}{4 D^2}.
\end{equation}
For a $10\,M_\odot$ merger at $D=1$\,pc from a $R_{\mathrm{throat}}=10^3$\,m
throat:
\begin{align}
  |\delta M_{\mathrm{throat}}^{\mathrm{GW}}|
  &\approx
  \frac{0.05\times 10\times 2\times 10^{30}\times 10^6}
       {4\times (3\times 10^{16})^2}
  \;\mathrm{kg} \nonumber\\
  &\approx
  \frac{10^{36}}{4\times 10^{33}}
  \;\mathrm{kg}
  \approx 250\;\mathrm{kg}.
\end{align}
This is slightly more substantial---around 250\,kg, or roughly three minutes
of human population on Earth.  As a fractional mass change of the wormhole:
$|\delta M_{\mathrm{throat}}^{\mathrm{GW}}|/M_{\mathrm{throat}}\sim 10^{-28}$.
The effect on S2's orbital period is
\begin{equation}
  \frac{\delta P}{P}\sim\frac{\delta M_{\mathrm{throat}}}{M_{\mathrm{Sgr\,A*}}}
  \approx\frac{250}{8\times 10^6\times 2\times 10^{30}}
  \approx 1.5\times 10^{-35},
\end{equation}
which is $\sim 30$ orders of magnitude below current observational precision.

Figure~\ref{fig:memory-estimate} summarizes the hierarchy of effects.

\begin{figure*}[tb]
\centering
\includegraphics[width=0.95\textwidth]{fig_memory_timeline.pdf}
\caption{Timeline of the radiation memory effect on the wormhole throat.
A radiation burst (GW or EM) approaches from mouth $B$ (top left), transits
the throat, and disperses at mouth $A$ (top right).  The throat mass contribution
$M_{\mathrm{throat}}(t)$ (bottom panel) decreases permanently after the transit by
$\delta M_{\mathrm{throat}}=-E/c^2<0$.  This is the throat memory: a
permanent, monopole-type change of \emph{negative} sign.  Compare with the
Dai--Stojkovic prediction (dashed) of a positive monopole growth: same physical
intuition, opposite sign, second-order magnitude.}
\label{fig:memory-timeline}
\end{figure*}

\subsection{Why only gravitational radiation can drive the nonlinear gap}

The single-ended gravity case has a residual open question at full nonlinear
order: could the gravitational field self-coupling (the Landau--Lifshitz
pseudo-tensor) change $Q_{\mathrm{wh}}^{\mathrm{grav}}$ even when matter stays
on one side?  The answer is that only gravitational radiation (not
electromagnetic radiation, not a static source) carries enough monopole-free
stress-energy to contribute through the throat, and even then:

\begin{enumerate}
  \item \emph{Electromagnetic radiation} is $T=0$ at the throat (no EM
    stress-energy at the throat location when EM radiation is not present
    there).  The gravitational backreaction of EM radiation passing through the
    throat is $O(G S/c^4)$ where $S$ is the Poynting flux, and gives the same
    throat deficit of negative sign computed above.

  \item \emph{Gravitational waves} are transverse-traceless (TT) at linear
    order.  The TT character means zero net monopole flux through the throat at
    $O(h)$; at $O(h^2)$, the Isaacson stress-energy has a non-zero
    energy density that can in principle drive a monopole change.  But the
    $O(h^2)$ contribution is the throat memory computed above---negative sign,
    second order.

  \item \emph{A static source on one side} produces a static field threading
    the throat.  By $Q_{\mathrm{wh}}$ conservation (linearized: exact;
    nonlinear: the static field does not radiate, so no GW are produced at the
    throat and the nonlinear correction is zero), the throat flux is frozen at
    its initial value.
\end{enumerate}

Therefore, across all physical processes, $Q_{\mathrm{wh}}^{\mathrm{grav}}$ either
(a) is conserved to all relevant orders, or (b) decreases at second order via
the radiation memory mechanism.  In no case does a first-order positive
monopole induction occur.

%% ============================================================
\section{One-Ended vs.\ Two-Ended: Comparison}
\label{sec:topology-comparison}
%% ============================================================

Table~\ref{tab:comparison} summarizes the conservation laws and D-S monopole
status for the two topologies.

\begin{table*}[t]
\centering
\caption{Conservation laws and Dai--Stojkovic monopole status for both wormhole topologies.}
\label{tab:comparison}
\begin{tabular}{lcccc}
\toprule
Topology & Conserved quantity & Lives in & Monopole induction & Mechanism \\
\midrule
Two-ended ($\mathbb{R}\times S^2$) & $M_\pm^{\mathrm{ADM}}$ each & $\partial\Sigma_\pm$ & Forbidden (Thm.~2) & HDA+RT \\
Two-ended & $Q_\pm$ each & $\partial\Sigma_\pm$ & Forbidden (Thm.~1) & Amp\`ere--Maxwell \\
Single-ended ($\mathbb{R}^3$+handle) & $Q_{\mathrm{wh}}$ & $H^1(\Sigma;\mathbb{R})$ & Forbidden (Thm.~3) & 4D Stokes \\
Single-ended & $Q_{\mathrm{wh}}^{\mathrm{grav}}$ & $H^1(\Sigma;\mathbb{R})$ & Forbidden at $O(h)$ & Linearized Bianchi \\
Both (radiation transit) & $M^{\mathrm{ADM}}$ & $\partial\Sigma$ & $-E/c^2$ at $O(h^2)$ & Memory (throat deficit) \\
\bottomrule
\end{tabular}
\end{table*}

The two-ended case has stronger conservation: \emph{separate} ADM masses at
each end, each individually conserved.  For a two-ended wormhole with both
mouths in the same universe (the D-S mathematical setup), this means that even
if matter crosses the throat from end $-$ to end $+$, $M_+^{\mathrm{ADM}}$
does not increase.  The throat flux $Q_{\mathrm{wh}}^{\mathrm{grav}}$ decreases
by the same amount to compensate.

The single-ended case has only total $M^{\mathrm{ADM}}$ conserved (trivially),
but the per-mouth flux $Q_{\mathrm{wh}}$ is separately conserved---and it is
this quantity that controls the apparent monopole at each mouth.  A source
staying on the $B$ side cannot change $Q_{\mathrm{wh}}$; a source transiting
does change $Q_{\mathrm{wh}}$ but only because it is now physically present at
mouth $A$.

In both cases the D-S observational signal---a time-varying far-zone monopole
driven by source motion on the other side---is absent.

%% ============================================================
\section{Multi-End Generalization}
\label{sec:multiend}
%% ============================================================

The extension to $N$ asymptotically flat ends is conceptually straightforward
and technically direct.

\begin{theorem}[Per-end conservation, $N$-ended]
\label{thm:N-ended}
Let $\Sigma$ be a smooth oriented asymptotically flat $N$-ended 3-manifold
satisfying the standard RT fall-off conditions at each end.  Assume (i) no
additional physical boundary components beyond the $N$ asymptotic ends, and
(ii) the existence of an admissible end-selecting lapse $N_a$ for each end
$E_a$ ($a=1,\ldots,N$).  Then:
\begin{enumerate}
  \item (Maxwell) $dQ_a/dt=0$ for each $a$ separately.
  \item (Gravity, conditional on admissibility) $dM_a^{\mathrm{ADM}}/dt=0$
    for each $a$ separately.
\end{enumerate}
\end{theorem}

\begin{proof}
The Maxwell case is identical to Theorem~\ref{thm:Q-conservation-2e} applied
to each end independently; the proof uses only the Amp\`ere--Maxwell equation
and the divergence theorem in the region exterior to $K$ at end $E_a$, and
does not depend on how many other ends exist.

The gravity case follows from exactly the same HDA+RT steps as the proof of
Theorem~\ref{thm:M-conservation}: the HDA identity \eqref{eq:HDA} holds for
any two end-selecting lapses $N_a$, $N_b$; their gradients vanish at all ends;
the resulting shift vector $\beta_{ab}^i$ is asymptotically trivial at all $N$
ends; and the weakly vanishing of $H[0,\beta_{ab}]$ follows from the same
argument, now with $N$ boundary terms all vanishing individually.
\end{proof}

\begin{remark}[Cohomological complexity]
For an $N$-ended spacetime with a connected handle structure, $H^2(\Sigma;\mathbb{R})$
can be higher-dimensional: an $N$-ended wormhole has up to $N-1$ independent
harmonic-flux parameters (one per independent non-contractible 2-cycle).  Each
such parameter is independently frozen, and per-end conservation holds for all
$N$ ends.  The D-S argument extends to any number of ends: none of them can
develop a time-varying monopole from internal source motion.
\end{remark}

\begin{remark}[Conditions and caveats]
The admissibility of end-selecting lapse functions (condition ii) becomes more
restrictive as $N$ grows.  For two ends, Regge and Teitelboim give an explicit
construction.  For $N\ge 3$, the lapse functions with different limiting values
at different ends must each satisfy the RT parity conditions---which requires
the geometry to be well-behaved at all $N$ ends simultaneously.  This is
expected to hold for the class of spacetimes considered in the wormhole
literature (asymptotically Schwarzschild at each end) but we flag it as a
condition rather than a theorem for general $N$.
\end{remark}

%% ============================================================
\section{What This Rules Out}
\label{sec:ruled-out}
%% ============================================================

\paragraph{Monopole induction.}
A source on end $E_-$ cannot produce a time-varying $Q_+$ or $M_+^{\mathrm{ADM}}$.
The opposite-end monopole is rigid.

\paragraph{Transit-driven monopole transfer.}
A body traversing the throat does not transfer monopole charge or ADM mass
between the two ends.  The asymptotic monopole data at each end remain at their
initial values throughout.

\paragraph{Far-zone monopole signal from hidden perturbers.}
Proposals to detect a wormhole at Sgr~A* by looking for a time-varying $1/r^2$
acceleration on S2~\cite{DaiStojkovic2019Observing,SimonettiEtAl2021SensitiveSearches}
require a varying opposite-end ADM monopole.  Theorems~\ref{thm:M-conservation}
and \ref{thm:Q-conservation-2e} forbid that signal for the two-ended topology;
Theorem~\ref{thm:Q-conservation-1e} forbids the analogous induction in the
single-ended topology.

%% ============================================================
\section{What This Does Not Rule Out}
\label{sec:not-ruled-out}
%% ============================================================

\paragraph{Higher multipoles.}
The conservation theorems pin only the $\ell=0$ asymptotic data.  Source motion
at end $E_+$ can perfectly well induce time-varying $\ell\ge 1$ multipoles at
end $E_-$.  These are exponentially suppressed in the long-throat regime
(Sec.~\ref{sec:attenuation}) and fall off as $r^{-(\ell+1)}$ rather than
$r^{-1}$, but they are real effects.

\paragraph{The harmonic-flux sector itself.}
A wormhole can carry a non-zero $Q_{\mathrm{wh}}$.  This is the genuine
Wheelerian charge-without-charge.  Per-end conservation does not deny its
existence; it denies that it is dynamically generated by local source transport.

\paragraph{Cross-throat influence at finite distance.}
Near-field effects---induction of higher multipoles, modification of the field
structure near the opposite mouth---are real.  The result rules out only the
leading $1/r$ (monopole) contribution at the far asymptotic end.

\paragraph{Quasi-local mass variation.}
Quasi-local mass functions evaluated on a finite sphere near one mouth can and
do respond to source motion on the other side, through the multipole tail
(Sec.~\ref{sec:quasilocal}).

\paragraph{Radiation transit effects.}
The throat memory effect of Sec.~\ref{sec:memory} is real.  It is second-order,
sign-reversed, and undetectable, but it is not forbidden.

%% ============================================================
\section{Multipole Attenuation Through the Throat}
\label{sec:attenuation}
%% ============================================================

The conservation argument removes the monopole signal.  This section estimates
how strongly the surviving higher-multipole channels are suppressed through a
long thin throat.

Model the throat as a uniform cylinder $S^2\times[0,L]$ of constant radius
$R$.  A static field $\Phi$ in the throat satisfies Laplace's equation:
\begin{equation}
  \Phi(z,\theta,\varphi)
  = \sum_{\ell m}f_{\ell m}(z)\,Y_{\ell m}(\theta,\varphi),
  \qquad
  \frac{d^2 f_{\ell m}}{dz^2} - \frac{\ell(\ell+1)}{R^2}f_{\ell m}=0.
\end{equation}

\begin{proposition}[Evanescent decay rate]
\label{prop:attenuation}
Inside the uniform cylindrical throat, each mode $\ell\ge 1$ is evanescent
with decay rate $k_\ell=\sqrt{\ell(\ell+1)}/R$ and transmission bounded above
by $e^{-\sqrt{\ell(\ell+1)}\,L/R}$.  The $\ell=0$ monopole has $k_0=0$ and
is not evanescent.
\end{proposition}

\begin{corollary}[Forbidden channel only is unsuppressed]
The unique multipole channel whose transmission is not exponentially suppressed
in $L/R$ is the $\ell=0$ monopole---which is forbidden by per-end conservation.
Every other channel is exponentially bounded.
\end{corollary}

For concreteness: the dipole bound is $e^{-\sqrt{2}L/R}\approx e^{-1.41L/R}$.
At aspect ratio $L/R=5$ the dipole is suppressed to $\sim 10^{-3}$; at $L/R=10$
to $\sim 10^{-6}$.  The $10^{-6}$ precision target of
\cite{DaiStojkovic2019Observing} for Sgr~A* would require $L/R\lesssim 10$
even for the dipole, and higher multipoles are much more strongly suppressed.
Figure~\ref{fig:attenuation} displays the evanescent bound as a function of
$L/R$ for the first four multipole channels.

\begin{figure}[t]
\centering
\includegraphics[width=0.85\textwidth]{fig3_attenuation.pdf}
\caption{Mode transmission through a uniform cylindrical wormhole throat of
length $L$ and radius $R$ versus aspect ratio $L/R$.  The $\ell=0$ monopole
(thick black) passes unattenuated but is forbidden by per-end conservation.
Every other channel is exponentially suppressed.  The shaded band marks a
$10^{-3}$ detection floor; for $L/R\gtrsim 5$ even the dipole drops below it.}
\label{fig:attenuation}
\end{figure}

%% ============================================================
\section{Phenomenological Implications}
\label{sec:phenomenology}
%% ============================================================

\subsection{The Sgr~A* proposal}

The Dai--Stojkovic Sgr~A* proposal~\cite{DaiStojkovic2019Observing} estimates
a time-varying acceleration $\Delta a\sim\mu R/r_p^3$ on S2, sourced by the
position-dependent piece $\mu R/r_p$ of the opposite-end ADM monopole (their
Eq.~37).  By Theorem~\ref{thm:M-conservation}, $M_-^{\mathrm{ADM}}$ at our
side does not vary.  The proposed signal is absent at the level of principle.

The Simonetti et al.\ extension~\cite{SimonettiEtAl2021SensitiveSearches}
makes the same monopole assumption explicit.  The conservation law applies
equally.

\subsection{Surviving signals}

What \emph{is} consistent with our results is:
\begin{enumerate}
  \item Higher-multipole cross-throat signals ($\ell\ge 1$), exponentially
    suppressed in $L/R$ and falling off as $r^{-(\ell+2)}$ at large $r$.  The
    signal-to-noise is far less favorable than the monopole estimates in the
    literature.
  \item The throat memory effect of Sec.~\ref{sec:memory}: a permanent but
    tiny negative-sign change in the effective throat mass, of order
    $E_{\mathrm{transit}}/c^2$.
\end{enumerate}

\subsection{Signal hierarchy}

Figure~\ref{fig:signal-hierarchy} compares the signal amplitudes for the three
main effects in the Sgr~A* scenario with S2 at orbital distance $r_{S2}\approx
1000$\,AU from Sgr~A*.

\begin{figure}[t]
\centering
\includegraphics[width=0.85\textwidth]{fig_signal_hierarchy.pdf}
\caption{Signal hierarchy for the Sgr~A* wormhole scenario.  The
Dai--Stojkovic claimed monopole signal (ruled out) sits far above any surviving
effect.  The dipole transmission through a $L/R=5$ throat is suppressed by
$\sim 10^{-3}$ relative to the nominal monopole amplitude.  The gravitational
wave throat memory from a $10\,M_\odot$ merger at 1\,pc is $\sim 30$ orders of
magnitude below current sensitivity.  The electromagnetic radiation throat
memory from a nearby star is further suppressed by another $\sim 10$ orders of
magnitude.}
\label{fig:signal-hierarchy}
\end{figure}

Any observational programme built on a time-varying far-zone monopole must
either (i) identify an alternative target in the higher-multipole channels, with
correspondingly less favourable signal-to-noise, or (ii) invoke physics outside
the hypotheses---for example, a time-varying harmonic sector sustained by
mechanisms not describable by smooth Einstein--Maxwell evolution with
compactly supported sources.

%% ============================================================
\section{Quasi-Local versus Asymptotic Mass}
\label{sec:quasilocal}
%% ============================================================

A natural objection: if I sit at finite radius $r_2$ on side 2 and measure the
gravitational pull of the wormhole mouth, surely that pull changes as the source
on side 1 moves?

Yes.  Quasi-local mass functions---Hawking, Brown--York, Misner--Sharp---each
defined on a finite 2-sphere, are sensitive to the full multipole structure of
the field and can respond to source motion on the opposite side.  What
Theorem~\ref{thm:M-conservation} pins is specifically $M^{\mathrm{ADM}}$, the
$r\to\infty$ limit.\cite{Szabados2009Quasilocal}

Concretely, the gravitational potential at finite $r_2$ on side 2 is
\begin{equation}
  \Phi(r_2) = -\frac{G M_-^{\mathrm{ADM}}}{r_2}
  + \sum_{\ell\ge 1}\frac{a_\ell(t)}{r_2^{\ell+1}},
\end{equation}
where the first term is rigid and the $a_\ell(t)$ multipole coefficients can
vary.  Attributing the full instantaneous $r_2^2\partial_{r_2}\Phi$ to an
``effective monopole charge'' is the technical error that produces the contrary
intuition.  The error vanishes only in the strict $r_2\to\infty$ limit.

%% ============================================================
\section{Discussion}
\label{sec:discussion}
%% ============================================================

The principal result of this paper is the per-end conservation of asymptotic
charge and ADM mass in the multi-ended asymptotically flat setting, and its
application to rule out---at the level of principle---the time-varying
opposite-end monopole on which several recent wormhole observational proposals
depend.  The result holds for both two-ended and single-ended topologies; the
conservation law lives in $H^2$ for the two-ended case and in $H^1$ for the
single-ended case, but in both cases it is exact at linear order and forbids
the claimed induction.

The secondary result is the identification of the correct subleading effect:
radiation transiting the throat imprints a permanent negative-sign monopole
change on the throat geometry, the throat memory effect.  This is a genuine
second-order contribution that partially vindicates the Dai--Stojkovic
intuition (something does happen when radiation crosses the throat) while
differing from their prediction in sign and being many orders of magnitude too
small to detect.

\paragraph{What remains genuinely interesting.}
The harmonic-flux sector $Q_{\mathrm{wh}}$ is the authentic Wheelerian
charge-without-charge: a topological flux supporting a Coulombic field at all
ends without a localized source.  A wormhole with a nonzero $Q_{\mathrm{wh}}$
either inherited it as primordial data or accumulated it via prior throat
transit.  The two histories are observationally distinct in principle (the
primordial case has no transit history; the accumulated case has a record of
prior crossings), but their asymptotic fields are identical.  The throat memory
effect---cumulative over all prior radiation transits---represents the only
known mechanism by which the throat geometry records a causal history of past
transits, albeit at an unobservably small level.

\paragraph{Open questions.}
The main open edge is the full nonlinear gravity treatment in the single-ended
case at order $h^2$.  The analysis above shows that the only physical process
capable of driving $Q_{\mathrm{wh}}^{\mathrm{grav}}$ at nonlinear order is
gravitational radiation transiting the throat, and that the resulting change is
of negative sign and second order.  A rigorous proof that this is the \emph{only}
nonlinear contribution---ruling out, e.g., gravitational wave scattering off
the throat that does not fully transit---would close the remaining gap.  We
expect this to follow from a careful analysis of the Israel junction conditions
on the throat worldtube (Appendix~\ref{app:israel}) extended to second order,
but leave the detailed proof to future work.

\paragraph{Scope.}
The paper's result is a sharpening of the no-go picture: it does not deny all
phenomenological interest in traversable wormholes.  Wormholes can have
nontrivial asymptotic charges; they can modulate higher multipoles across the
throat; they accumulate a throat memory from radiation transit.  What they
cannot do, under the hypotheses stated here, is modulate the far-zone monopole
by ordinary interior source motion.  A legitimate observational programme would
have to target higher-multipole signals, with their faster radial fall-off and
exponential evanescent suppression in the long-throat limit, or the tiny
second-order throat memory signature---both far more challenging than the
monopole signal that has been claimed.

%% ============================================================
\appendix
%% ============================================================

\section{Cohomological Characterisation of the Harmonic Sector}
\label{app:cohomology}

We collect the standard cohomological facts used in the body of the paper.

On a spacelike slice $\Sigma$, the Maxwell constraint $d(\star F)=4\pi j$
implies $d(\star F)=0$ in source-free regions, so $\star F$ is closed there.
On $\Sigma\cong\mathbb{R}\times S^2$, closed need not mean exact: the single
non-exact direction in $H^2(\Sigma;\mathbb{R})\cong\mathbb{R}$ is the
Wheelerian harmonic-flux sector.  A normalised generator is
\begin{equation}
  h = \frac{1}{4\pi}\Omega,
\end{equation}
where $\Omega$ is the standard area 2-form on $S^2$, pulled back via the
projection $\mathbb{R}\times S^2\to S^2$.

$Q_{\mathrm{wh}}$ is conserved whenever the current support, over the time
interval $[0,T]$, is disjoint from some representative 2-cycle $S$ of the
$H_2$ generator.  The 4-dimensional Stokes theorem on $S\times[0,T]$ then
gives:
\begin{equation}
  \int_{S\times\{T\}}\star F - \int_{S\times\{0\}}\star F
  = 4\pi\int_{S\times[0,T]}\star J = 0.
\end{equation}
If the current traverses every representative of $H_2$ (i.e., crosses the
throat), $Q_{\mathrm{wh}}$ shifts by the transported charge, while $Q_\pm$
remain fixed.

\section{Explicit Transit Calculation in the Short-Throat Model}
\label{app:transit}

We make the harmonic-flux accounting explicit in the Dai--Stojkovic
cut-and-paste geometry.  For a source of total charge $q$ at position $A$ on
side $+$, the Dai--Stojkovic matching conditions select
\begin{equation}
Q_-^{\mathrm{DS}}(A) = \frac{qR}{2A}.
\end{equation}
Any actual Maxwell evolution starting from source at $A_0$ gives
$Q_\pm(t)=Q_\pm(0)$ throughout (Theorem~\ref{thm:Q-conservation-2e}).  The
end-state at $A_1$ has $Q_-=qR/(2A_0)$, not $qR/(2A_1)$.  The required
harmonic-flux drift is:
\begin{equation}
\Delta Q_{\mathrm{wh}} = \frac{qR}{2}\!\left(\frac{1}{A_0}-\frac{1}{A_1}\right),
\end{equation}
which must be zero by cohomological rigidity for in-end source motion.

For throat transit ($A$ passing through $R$), $Q_\pm$ stay at their initial
values; $Q_{\mathrm{wh}}$ shifts by $q$; the configuration on side $+$ after
transit is a ``charge without charge'' of strength $q$ with no localized
source on that side.

\section{Throat Memory: Israel Junction Condition Calculation}
\label{app:israel}

We make the throat memory calculation explicit for a thin-shell traversable
wormhole model~\cite{Visser1995LorentzianWormholes}.

\paragraph{Background.}
The wormhole throat is modeled as a spherically symmetric 2-shell of radius
$R_0$ with surface energy density $\sigma_0<0$ (exotic matter, supporting the
throat against collapse).  The Israel junction conditions
relate the jump in extrinsic curvature $[K_{ab}]$ to the shell stress-energy
$S_{ab}=-\sigma_0 h_{ab}+\tau_{ab}$, where $\tau_{ab}$ is the tangential
pressure contribution.  The effective Schwarzschild mass of the shell is
$M_{\mathrm{throat}}=4\pi R_0^2 \sigma_0<0$.

\paragraph{Radiation transit.}
Consider a burst of radiation (electromagnetic or gravitational) of energy $E$
transiting the throat in a proper time $\delta\tau$.  The burst deposits
stress-energy $T_{\mu\nu}^{\mathrm{rad}}$ at the shell location.  To leading
order in $E/(M_{\mathrm{throat}}c^2)$, the junction conditions are perturbed as
\begin{equation}
  [K_{ab}]_{\mathrm{new}}
  = [K_{ab}]_0 - 8\pi G\,\Delta S_{ab},
\end{equation}
where $\Delta S_{ab}$ is the integrated impulse from the radiation.  For a
spherically symmetric burst, the $l=0$ component of $\Delta S_{ab}$ is
$\Delta\sigma = E/(4\pi R_0^2 c^2)$, giving
\begin{equation}
  \Delta M_{\mathrm{throat}} = 4\pi R_0^2\,\Delta\sigma = \frac{E}{c^2}.
\end{equation}
Wait---this gives a \emph{positive} change to the shell mass: the exotic matter
shell becomes less negative (less exotic) by the energy deposited.  However,
the \emph{gravitational flux contribution} to $M_+^{\mathrm{ADM}}$ from the
throat is $-M_{\mathrm{throat}}$ (since the throat is the exotic matter
holding the wormhole open; from the $+$ end, it appears as a negative mass
contribution).  Therefore:
\begin{equation}
  \delta(\text{throat contribution to }M_+^{\mathrm{ADM}})
  = -\Delta M_{\mathrm{throat}} = -\frac{E}{c^2}.
\end{equation}
This is the throat mass deficit: the throat contributes $-E/c^2$ less to
$M_+^{\mathrm{ADM}}$ after the transit.  Combined with the $+E/c^2$
contribution of the radiation now on the $A$ side, $M_+^{\mathrm{ADM}}$
is exactly conserved, as required.

\paragraph{Permanent character.}
After the burst has passed and the shell has re-equilibrated, the shell mass
$M_{\mathrm{throat}}$ is permanently shifted by $+E/c^2$ (slightly less
exotic).  The effective gravitational flux contribution to $M_+^{\mathrm{ADM}}$
from the throat is permanently decreased by $E/c^2$.  This is the throat
memory.

\bibliographystyle{apsrev4-2}
\bibliography{references}

\end{document}
